\begin{abstract}
Bus timetable planning and real-time holding control are operationally coupled yet typically addressed in isolation. This paper presents TransitDuet, a bi-level reinforcement learning framework that jointly optimises rolling dispatch timetables and station-level holding decisions under stochastic demand and traffic conditions. The upper policy selects discrete planned-headway adjustments that define a dynamic rolling timetable, while the lower policy, a robust ensemble soft actor--critic controller with Lagrangian headway constraints, executes this timetable through station holding. Cross-level coordination operates through a dynamic timetable channel, holding-feedback state augmentation, and per-dispatch hindsight credit assignment, with fleet-size feasibility maintained by measurement-driven online gradient descent penalty adaptation. Evaluated in a calibrated bidirectional bus-corridor microsimulation with elastic fleet budgets across ten random seeds, TransitDuet achieves a mean composite cost of 1.634, which is 3.2\% lower than a target-headway-only variant (1.687) and 11.5\% lower than the best fixed-timetable baseline (1.846). Mean total passenger time is 18.15 minutes, compared with 19.30 and 19.92 minutes for these respective comparators. These results demonstrate that jointly optimising dispatch planning and holding control yields measurable performance gains over single-level holding controllers and decoupled configurations on the evaluated corridor.
\end{abstract}
